\PassOptionsToPackage{dvipsnames}{xcolor}

\documentclass[12pt]{beamer}
\usetheme{default}
\usecolortheme{crane}

\usepackage[T1]{fontenc}
\usepackage[slovak]{babel}

\usepackage{amsmath}
\usepackage{amsmath, amssymb}
\usepackage{hyperref, url}
\usepackage{graphicx}
\usepackage{array}
\usepackage{alltt}

%\setbeamersize{text margin left=1pt,text margin right=1pt}
\setbeamertemplate{footline}[frame number]
\beamertemplatenavigationsymbolsempty

% https://www.overleaf.com/learn/latex/Using_colours_in_LaTeX
\def\blue#1{\textcolor{Cerulean}{#1}}
\def\green#1{\textcolor{LimeGreen}{#1}}

% database-related stuff
\DeclareMathOperator{\join}{\bowtie}
\DeclareMathOperator{\antijoin}{\rhd}

\DeclareMathOperator{\lubi}{lubi}
\DeclareMathOperator{\capuje}{capuje}
\DeclareMathOperator{\navstivil}{navstivil}
\DeclareMathOperator{\vypil}{vypil}
\DeclareMathOperator{\answer}{answer}

\title{Rekurzia}
\author{Ján Mazák}
\institute{FMFI UK Bratislava}
\date{}


\begin{document}

\frame{\titlepage}

\begin{frame}[fragile]{Rekurzia: cez CTE (Common Table Expressions)}
\alert{WITH RECURSIVE} umožňuje rekurzívne sa odvolať na reláciu.
\bigskip
\begin{alltt}
\alert{WITH RECURSIVE} cte_name AS(
    SELECT ...  ---- non-recursive term
  UNION [ALL]
    SELECT ...  ---- recursive term
)
SELECT * FROM cte_name;
\end{alltt}
\end{frame}

\begin{frame}[fragile]{Rekurzia}
%\begin{minipage}{.4\pdfpagewidth}
\begin{alltt}
\alert{WITH RECURSIVE} t(n) AS (
    SELECT 1
  UNION ALL
    SELECT n+1 FROM t \green{WHERE n < 5}
)
SELECT n FROM t;
\end{alltt}
%\end{minipage}
\end{frame}

\begin{frame}[fragile]{Výpočet rekurzie --- iterácia (seminaive evaluation)}
Naivná evaluácia v každom kroku znova vyhodnotí rekurzívny člen nad všetkými doterajšími riadkami; seminaivná spracúva len riadky pridané v predchádzajúcom kroku.
\begin{enumerate}
\item Evaluate the \blue{non-recursive} term. For UNION (but not UNION ALL), discard duplicate rows.
Include all remaining rows in the result of the recursive query, and also place them in a temporary working table.

\item Evaluate the \blue{recursive} term, substituting the current contents of the working table for the recursive self-reference.
For UNION, discard duplicate rows and rows that duplicate any previous result row.
Include all remaining rows in the result of the recursive query, and also place them in a temporary intermediate table.

\item If the intermediate table is empty, stop. Otherwise replace the working table with the intermediate table and go to 2.
\end{enumerate}
\end{frame}

\begin{frame}[fragile]{Rekurzia}
Rekurzívny odkaz možno použiť len raz a pre jednu reláciu (PostgreSQL).
\alert{Nekorektné} využitie:
%\begin{minipage}{.4\pdfpagewidth}
\begin{alltt}
WITH RECURSIVE t(n) AS (
    SELECT 1
  UNION ALL
    (SELECT n+1 FROM \blue{t} WHERE n < 5
        UNION ALL
     SELECT n-1 FROM \blue{t} WHERE n < 5)
)
SELECT n FROM t;
\end{alltt}
%\end{minipage}
\end{frame}

\begin{frame}[fragile]{Rekurzia}
Rekurzívny výpočet nemusí skončiť, zodpovednosť je na autorovi dotazu.
\bigskip

\begin{alltt}
WITH RECURSIVE t(n) AS (
    SELECT 1
  UNION ALL
    SELECT n+1 FROM t
)
SELECT n FROM t \green{WHERE n < 5};
\end{alltt}
\bigskip
(v PostgreSQL sa takto dotaz nezastaví)
\end{frame}

\begin{frame}[fragile]{Rekurzia}
Rekurzia zvyšuje vyjadrovaciu silu jazyka SQL.
\begin{itemize}
\item Rekurziou sa dá vyjadriť \blue{tranzitívny uzáver},\\ t.\,j. existencia cesty neobmedzenej dĺžky.
\item Bez rekurzie možno vyjadriť len cestu fixnej dĺžky (pomocou opakovaného joinu).
\end{itemize}
\end{frame}

\begin{frame}[fragile]{Prehľadávanie grafov}
Pomocou rekurzie možno zapísať prehľadávanie stromu.\\
Bez určenia počiatočného vrcholu ide o prehľadávanie zo všetkých vrcholov naraz; počiatočný vrchol určíme podmienkou v prvej časti (WHERE t.vFrom = \dots).\\
\bigskip

\begin{alltt}
WITH RECURSIVE \alert{search_tree}(vFrom, vTo) AS (
    SELECT t.vFrom, t.vTo
    FROM tree t
  UNION ALL
    SELECT st.vFrom, t.vTo
    FROM \alert{search_tree} st, tree t
    WHERE t.vFrom = st.vTo
)
SELECT * FROM search_tree;
\end{alltt}
\end{frame}

\begin{frame}[fragile]{Prehľadávanie grafov}
\small
Poradie riadkov vo výsledku nie je zaručené, ale usporiadaním podľa hĺbky dostaneme vrcholy v poradí \alert{prehľadávania do šírky}.
\bigskip
\begin{alltt}
WITH RECURSIVE search_tree(vFrom, vTo, \alert{depth}) AS (
    SELECT t.vFrom, t.vTo, \alert{1}
    FROM tree t
  UNION ALL
    SELECT st.vFrom, t.vTo, \alert{depth + 1}
    FROM search_tree st, tree t
    WHERE t.vFrom = st.vTo
)
SELECT * FROM search_tree \alert{ORDER BY depth};
\end{alltt}
\end{frame}

\begin{frame}[fragile]{Prehľadávanie grafov}
Poradie riadkov vo výsledku nie je zaručené, ale usporiadaním podľa cesty dostaneme vrcholy v poradí \alert{prehľadávania do hĺbky}.
\bigskip
\begin{alltt}
WITH RECURSIVE search_tree(vFrom, vTo, \alert{path}) AS (
    SELECT t.vFrom, t.vTo, \alert{ARRAY[t.vFrom, t.vTo]}
    FROM tree t
  UNION ALL
    SELECT st.vFrom, t.vTo, \alert{path || t.vTo}
    FROM search_tree st, tree t
    WHERE t.vFrom = st.vTo
)
SELECT * FROM search_tree \alert{ORDER BY path};
\end{alltt}
\end{frame}

\begin{frame}[fragile]{Prehľadávanie grafov}
Ak chceme ustrážiť koniec rekurzie, treba sa vyhýbať \alert{zacykleniu}.
\bigskip
{\small
\begin{alltt}
WITH RECURSIVE
    search_graph(vFrom, vTo, depth, \green{path}, \blue{is_cycle})
AS (
    SELECT g.vFrom, g.vTo, 1,
        \green{ARRAY[g.vFrom, g.vTo]}, \blue{false}
    FROM graph g
  UNION ALL
    SELECT sg.vFrom, g.vTo, sg.depth + 1,
        \green{path || g.vTo}, \blue{g.vTo = ANY(path)}
    FROM search_graph sg, graph g
    WHERE g.vFrom = sg.vTo AND \blue{NOT is_cycle}
)
SELECT * FROM search_graph ORDER BY is_cycle, path;
\end{alltt}
}
\end{frame}

\begin{frame}[fragile]{Prehľadávanie grafov}
SQL po pridaní rekurzie je turingovsky úplný jazyk
(vyjadríte v ňom riešenie pre všetky veci programovateľné v Pythone či C).
Náčrt dôkazu a ukážka riešenia zložitých problémov v SQL, napr. ASCII Mandelbrotovej množiny:
\href{https://stackoverflow.com/questions/900055/is-sql-or-even-tsql-turing-complete}{link}
\\[5mm]
\pause

Neznamená to však, že túto silu musíte aj využiť.
Napr. pri prehľadávaní grafu treba zvážiť, či nie je jednoduchšie vypýtať si celý graf
z databázy a prehľadávanie zapísať ako program vo vhodnejšom jazyku.
Podobne pri iných komplikovaných problémoch (napr. riešených backtrackingom).\\[5mm]

Hierarchické dáta sa však v databázach občas vyskytujú, treba preto mať aspoň predstavu, ako sa s nimi vysporiadať.
\end{frame}


\begin{frame}{Literatúra}
\begin{itemize}
\item {\scriptsize\url{https://www.postgresql.org/docs/current/queries-with.html}}
\item {\scriptsize\url{https://www.sqlite.org/lang_with.html}}
\end{itemize}
\end{frame}


\begin{frame}{Na zamyslenie}
\begin{itemize}
\item Dá sa každá rekurzia zapísať pomocou iterácie?\\ A čo iterácia pomocou rekurzie?
\item Vysvetlite rozdiel medzi naivnou a seminaivnou evaluáciou rekurzívneho dotazu.
\item Vymyslite otázku, ktorá sa nedá vyjadriť v SQL ani po pridaní rekurzie.
    (Bežne sa pýtame na záznamy, to sa vyjadriť dá, ale skúste sa pýtať na niečo iné.)
\end{itemize}
\end{frame}


\begin{frame}[fragile]{Úlohy}
Čo počíta tento dotaz?
\bigskip
{\small
\begin{alltt}
WITH RECURSIVE t(s) AS (
    SELECT 'foo'
  UNION ALL
    SELECT s || 'bar' FROM t WHERE length(s) < 20
)
SELECT s FROM t ORDER BY s;
\end{alltt}
}
\bigskip
\pause
Generuje reťazce v tvare \verb|foo(bar)|$^*$ postupne od najkratších.
\end{frame}


\begin{frame}[fragile]{Úlohy}
Súčet čísel od 1 do 100
\pause
\bigskip
{\small
\begin{alltt}
WITH RECURSIVE t(n) AS (
    VALUES (1)
  UNION ALL
    SELECT n+1 FROM t WHERE n < 100
)
SELECT sum(n) FROM t;
\end{alltt}
}
\end{frame}


\begin{frame}[fragile]{Úlohy}
Čo počíta tento dotaz?\\
Zmení sa to, ak UNION nahradíme UNION ALL?
\bigskip
{\small
\begin{alltt}
WITH RECURSIVE t(n) AS (
    SELECT 1
  \blue{UNION}
    SELECT 10-n FROM t
)
SELECT * FROM t ORDER BY n;
\end{alltt}
}
\end{frame}


\begin{frame}[fragile]{Úlohy}
Skončí výpočet tohto dotazu?
\bigskip
{\small
\begin{alltt}
WITH RECURSIVE t(n) AS (
    VALUES (1)
  UNION ALL
    SELECT n+1 FROM t
)
SELECT * FROM t LIMIT 10;
\end{alltt}
}
\bigskip
\pause
V PostgreSQL áno, ale vo všeobecnosti to závisí od implementácie LIMIT a rekurzie.
\end{frame}



\end{document}
